\documentclass{article}
\usepackage{pythontex}
\usepackage{amsmath,amssymb}
\usepackage{gensymb}

\begin{document}

\begin{pycode}
import random
import math

# Generate random values for the problem
angle = random.randint(30, 60)
initial_velocity = random.randint(20, 30)
time = random.randint(2, 5)

# Calculate the horizontal and vertical components of velocity
v0x = initial_velocity * math.cos(math.radians(angle))
v0y = initial_velocity * math.sin(math.radians(angle))

# Calculate the maximum height and horizontal range
h_max = (v0y ** 2) / (2 * 9.8)
range = v0x * time

\end{pycode}

\section*{Projectile Motion Problem}

A projectile is launched with an initial velocity of \(\py{initial_velocity}\, \text{m/s}\) at an angle of \(\py{angle}$\degree$\) with the horizontal. Calculate the following:

\begin{enumerate}
    \item The initial horizontal component of velocity (\(v_{0x}\)).
    \item The initial vertical component of velocity (\(v_{0y}\)).
    \item The maximum height reached by the projectile (\(h_{\text{max}}\)).
    \item The horizontal range of the projectile (\(R\)).
\end{enumerate}

\subsection*{Solution}

\begin{enumerate}
    \item The initial horizontal component of velocity (\(v_{0x}\)) can be calculated using the formula:
    \[v_{0x} = v_0 \cos(\theta)\]
    Substituting the given values, we have:
    \[v_{0x} = \py{round(v0x,2)}\, \text{m/s}\]

    \item The initial vertical component of velocity (\(v_{0y}\)) can be calculated using the formula:
    \[v_{0y} = v_0 \sin(\theta)\]
    Substituting the given values, we have:
    \[v_{0y} = \py{round(v0y,2)}\, \text{m/s}\]

    \item The maximum height reached by the projectile (\(h_{\text{max}}\)) can be calculated using the formula:
    \[h_{\text{max}} = \frac{{v_{0y}^2}}{{2g}}\]
    Substituting the given values, we have:
    \[h_{\text{max}} = \py{round(h_max,2)}\, \text{m}\]

    \item The horizontal range of the projectile (\(R\)) can be calculated using the formula:
    \[R = v_{0x} \cdot t\]
    Substituting the given values, we have:
    \[R = \py{round(range,2)}\, \text{m}\]
\end{enumerate}

    

\end{document}
